% =====================================================================
%  MSc Thesis — Imperial College London, Mathematics and Finance
%  Arbitrage-Free Implied Volatility Surface Construction:
%  From Parametric Models to Neural Operators
%  Structure follows MSc_thesis_Template_2026 (title page, declaration,
%  acknowledgements, abstract, ToC, unnumbered introduction, chapters,
%  appendix, bibliography).
% =====================================================================
\documentclass[11pt,a4paper,twoside]{report}

% ---------- packages ----------
\usepackage[utf8]{inputenc}
\usepackage[T1]{fontenc}

\usepackage{amsmath,amssymb,amsthm,mathtools}
\usepackage[margin=2.5cm]{geometry}
\usepackage{graphicx}
\usepackage{subcaption}
\usepackage{booktabs}          % professional tables
\usepackage{siunitx}           % aligned numeric columns
\usepackage{xcolor}
\usepackage[colorlinks=true, linkcolor=blue!60!black,
            citecolor=blue!60!black, urlcolor=blue!60!black]{hyperref}
\usepackage[capitalise,noabbrev]{cleveref}
\usepackage{listings}
\usepackage[expansion=false]{microtype}
\usepackage{emptypage}

% ---------- theorem environments (template style) ----------
\newtheorem{theorem}{Theorem}[section]
\newtheorem{lemma}[theorem]{Lemma}
\newtheorem{proposition}[theorem]{Proposition}
\newtheorem{corollary}[theorem]{Corollary}
\theoremstyle{definition}
\newtheorem{definition}[theorem]{Definition}
\newtheorem{example}[theorem]{Example}
\newtheorem{assumption}[theorem]{Assumption}
\theoremstyle{remark}
\newtheorem{remark}[theorem]{Remark}

% ---------- notation macros ----------
\newcommand{\E}{\mathbb{E}}
\newcommand{\R}{\mathbb{R}}
\newcommand{\Q}{\mathbb{Q}}
\newcommand{\iv}{\sigma_{\mathrm{IV}}}
\newcommand{\dd}{\mathrm{d}}
\DeclareMathOperator*{\argmin}{arg\,min}
\newcommand{\relu}{\mathrm{ReLU}}

% ---------- code listings ----------
\lstset{basicstyle=\ttfamily\small, breaklines=true, frame=single,
        language=Python, showstringspaces=false,
        keywordstyle=\color{blue!70!black}, commentstyle=\color{green!50!black}}

% ---------- metadata: EDIT THESE ----------
\newcommand{\thesisauthor}{Hamza \textsc{LASTNAME}}   % <-- EDIT
\newcommand{\thesiscid}{01234567}                      % <-- EDIT
\newcommand{\thesistitle}{Arbitrage-Free Implied Volatility Surface Construction:\\ From Parametric Models to Neural Operators}
\newcommand{\thesisyear}{2025--2026}


% ---------- compile-safe figures & todo notes ----------
\newcommand{\TODO}[1]{\textcolor{red}{[TODO: #1]}}
\newcommand{\maybegraphics}[2][]{\IfFileExists{#2}{\includegraphics[#1]{#2}}{\fbox{\parbox[c][5cm][c]{0.85\linewidth}{\centering \texttt{\detokenize{#2}}\\[4pt] \small (export this figure from the notebook and place it here)}}}}

\begin{document}

% =====================  TITLE PAGE  =====================
\begin{titlepage}
\begin{center}
  {\Large Imperial College London}\\[2pt]
  {\large Department of Mathematics}\\[3.2cm]
  {\LARGE\bfseries \thesistitle}\\[2.6cm]
  {\large Author: \thesisauthor{} (CID: \thesiscid)}\\[3.2cm]
  {\large A thesis submitted for the degree of}\\[6pt]
  {\large\itshape MSc in Mathematics and Finance, \thesisyear}
\end{center}
\end{titlepage}

% =====================  DECLARATION  =====================
\chapter*{Declaration}
The work contained in this thesis is my own work unless otherwise stated.
\cleardoublepage

% =====================  ACKNOWLEDGEMENTS  =====================
\chapter*{Acknowledgements}
% EDIT: thank supervisor, family, etc.
I would like to thank my supervisor for their guidance throughout this project.
\cleardoublepage

% =====================  ABSTRACT  =====================
\chapter*{Abstract}
The implied volatility surface is the central object of equity derivatives markets, yet it is only ever observed as a sparse, noisy and irregular cloud of option quotes. Reconstructing a complete, smooth and arbitrage-free surface from these observations is an ill-posed inverse problem, and the quality of the reconstruction directly affects pricing, hedging and risk management. This thesis studies the problem across the full methodological spectrum, on a single unified S\&P~500 dataset built from OptionMetrics. We first replicate the parametric benchmark of Gatheral and Jacquier---the SVI and SSVI families, their no-arbitrage theory and calibration recipes---validating every component against the original paper. We then implement a derivative-constrained deep smoother in the spirit of Ackerer et al.\ and Hoshisashi et al., in which a neural corrector multiplies an SSVI prior and no-arbitrage conditions are enforced through automatic differentiation, and finally a neural operator following the operator deep smoothing approach, which learns the smoothing map itself and applies to unseen days without retraining. All methods are evaluated under a common protocol: hold-out reconstruction error, a systematic static-arbitrage audit based on the Durrleman condition, robustness to quote sparsity, and behaviour across market regimes including the COVID-19 crash. As a contribution, we quantify the low-data regime of operator smoothing available to academic research and measure whether pre-training on synthetic SSVI surfaces closes the gap to the data-rich industrial setting.
\cleardoublepage

% =====================  FRONT MATTER  =====================
\tableofcontents
\listoffigures
\listoftables
\cleardoublepage

% =====================  CHAPTERS  =====================
\chapter*{Introduction}
\addcontentsline{toc}{chapter}{Introduction}

When an equity option is priced with the Black--Scholes formula, a single parameter is not directly observable: the volatility $\sigma$. The \emph{implied volatility} of a quoted option is the value of $\sigma$ that reproduces its market price, and when computed across strikes $K$ and expiries $\tau$ it traces out the \emph{implied volatility surface} $(k,\tau)\mapsto \iv(k,\tau)$, with $k=\log(K/F)$ the log-moneyness against the forward. The surface is the market's summary of the risk-neutral distribution of the underlying at every horizon; it is the input to exotic pricing, hedging and risk systems, and its shape---the skew, the term structure---is economically informative in its own right.

The market, however, never exhibits a surface. It exhibits a \emph{sparse, noisy, irregular} cloud of quotes: only certain strikes and expiries trade, liquidity concentrates near the money, and bid--ask spreads widen rapidly in the wings. Reconstructing a complete, smooth surface from these observations is an ill-posed inverse problem: infinitely many surfaces interpolate the same quotes. The discipline that selects an admissible one is \emph{absence of static arbitrage}: the reconstructed prices must not allow riskless profit through butterfly or calendar strategies. These constraints admit a precise analytical form---non-negativity of the Durrleman function $g$ for butterfly arbitrage, monotonicity of total variance in maturity for calendar arbitrage---which makes them both testable and, crucially for this thesis, \emph{trainable}.

Three generations of methods address the problem. Parametric families, culminating in SVI and SSVI \cite{gatheral2014arbitrage}, offer interpretability and provable no-arbitrage guarantees at the cost of rigidity. Per-surface neural smoothers \cite{ackerer2020deep, chataigner2020deep, hoshisashi2024noarbitrage} recover flexibility by penalising arbitrage through the network's automatically-differentiated derivatives, but require one model per surface. Neural operators \cite{gonon2024operator} learn the smoothing \emph{map} itself: trained once across many days, a single discretisation-invariant model smooths any quote set without retraining.

\paragraph{Contributions.} This thesis makes four contributions.
\begin{enumerate}
  \item \textbf{A unified benchmark.} All three generations are implemented and evaluated on the \emph{same} daily S\&P~500 dataset (OptionMetrics, 2018--2022), under the \emph{same} protocol: hold-out reconstruction error, a dense-grid static-arbitrage audit, robustness to quote sparsity, and per-regime analysis. Published papers each use their own data and metrics; to our knowledge no like-for-like comparison exists.
  \item \textbf{A validated replication.} The complete Gatheral--Jacquier toolkit---three SVI parameterisations and their exact conversions, the butterfly-elimination algorithm, the calibration recipes, the SSVI no-arbitrage theorems, arbitrage-free interpolation---is replicated and validated against the paper's own numerical examples, correcting in passing a typographical omission in the printed statement of the Jump-Wings inversion (Lemma~3.2).
  \item \textbf{A systematic arbitrage audit.} Beyond fit error, every method is scored on \emph{how much} static arbitrage it leaves on real data---violation rates and magnitudes of the Durrleman condition on dense grids, cross-slice crossedness, and term-structure inversions---quantities the literature rarely reports.
  \item \textbf{The low-data operator experiment.} Operator deep smoothing was demonstrated on $\sim$60 million intraday quotes; academic data access is daily, two orders of magnitude smaller. We measure operator performance in this regime and test whether pre-training on synthetic SSVI surfaces closes the gap (regimes R1/R2/R3: real-only, synthetic-only, pre-train-then-fine-tune).
\end{enumerate}

\paragraph{Outline.} Chapter~\ref{ch:foundations} sets the financial and mathematical foundations. Chapter~\ref{ch:literature} reviews the three methodological generations. Chapter~\ref{ch:data} documents the data pipeline. Chapters~\ref{ch:parametric}--\ref{ch:operator} develop the three method families, each validated on synthetic ground truth before being applied to market data. Chapter~\ref{ch:evaluation} performs the unified comparison, and the conclusion discusses limitations and extensions.

\chapter{Foundations}\label{ch:foundations}

\section{Options, Black--Scholes and implied volatility}
Let $S$ denote the underlying, $F=F(t,T)$ its forward for delivery at $T$, and $D=D(t,T)$ the discount factor. Under the Black (forward) model, the price of a European call of strike $K$ and time-to-maturity $\tau=T-t$ is
\begin{equation}\label{eq:black}
  C(K,\tau;\sigma) \;=\; D\left[\,F\,N(d_+) - K\,N(d_-)\,\right],
  \qquad d_{\pm} \;=\; \frac{\log(F/K)}{\sigma\sqrt{\tau}} \pm \frac{\sigma\sqrt{\tau}}{2},
\end{equation}
with $N$ the standard normal cdf. Since $\partial C/\partial\sigma>0$ (the vega), \cref{eq:black} can be inverted: the \emph{implied volatility} $\iv(K,\tau)$ of a quote $C^{\mathrm{mkt}}$ is the unique root of $C(K,\tau;\sigma)=C^{\mathrm{mkt}}$. The empirical fact that $\iv$ varies with $(K,\tau)$---the \emph{smile} and the term structure---contradicts the constant-$\sigma$ assumption and gives rise to the implied volatility surface.

\section{Coordinates: log-moneyness and total variance}
Two changes of variable make the no-arbitrage theory tractable. The \emph{log-moneyness} $k=\log(K/F)$ centres the smile on the forward. The \emph{total variance}
\begin{equation}
  w(k,\tau) \;=\; \iv^2(k,\tau)\,\tau
\end{equation}
is the natural vertical coordinate: both static no-arbitrage conditions take their simplest form in $w$.

\section{Static no-arbitrage}\label{sec:noarb}
Following \cite{gatheral2014arbitrage, roper2010arbitrage}, a surface is free of static arbitrage if and only if each maturity slice is free of butterfly arbitrage and the family of slices is free of calendar arbitrage.

\begin{definition}[Butterfly arbitrage and the Durrleman function]
For a slice $w(\cdot)$ at fixed maturity, define
\begin{equation}\label{eq:durrleman}
  g(k) \;=\; \Big(1-\frac{k\,w'(k)}{2\,w(k)}\Big)^{2} \;-\; \frac{w'(k)^{2}}{4}\Big(\frac{1}{w(k)}+\frac14\Big) \;+\; \frac{w''(k)}{2}.
\end{equation}
The slice is free of butterfly arbitrage if and only if $g(k)\ge 0$ for all $k$ and call prices tend to zero at extreme strikes (for SVI, the wing-slope condition $b(1+|\rho|)<2$).
\end{definition}

The financial content of \cref{eq:durrleman} is the Breeden--Litzenberger identity: the risk-neutral density is proportional to $\partial^2 C/\partial K^2$, and in total-variance coordinates
\begin{equation}
  p(k) \;=\; \frac{g(k)}{\sqrt{2\pi\,w(k)}}\,\exp\!\Big(-\tfrac12 d_-(k)^2\Big),
  \qquad d_-(k) = -\frac{k}{\sqrt{w(k)}}-\frac{\sqrt{w(k)}}{2},
\end{equation}
so $g<0$ is exactly a negative implied density.

\begin{definition}[Calendar arbitrage]
The surface is free of calendar arbitrage if and only if $\tau\mapsto w(k,\tau)$ is non-decreasing for every $k$; equivalently, total-variance slices never cross.
\end{definition}

\Cref{eq:durrleman} involves only $w$, $w'$ and $w''$. This is the property exploited throughout the thesis: for parametric models the derivatives are closed-form; for neural models they are obtained \emph{exactly} by automatic differentiation, turning the no-arbitrage conditions into differentiable training penalties.

% FIGURE (export from NB02, Section 2): Vogt smile three panels
\begin{figure}[t]
  \centering
  \maybegraphics[width=\textwidth]{figures/nb02_vogt_arbitrage.png}
  \caption[The Vogt smile: hidden butterfly arbitrage]{The Vogt SVI parameters of \cite[Ex.~3.1]{gatheral2014arbitrage}. The total-variance slice (left) looks innocent, yet the Durrleman function dips below zero (centre) and the implied risk-neutral density turns negative (right): a genuine butterfly arbitrage invisible to the naked eye. Figure produced by our replication (Notebook~02).}
  \label{fig:vogt}
\end{figure}

\section{The reconstruction problem}
Given a day's quotes $\{(k_i,\tau_i,{\iv}_i)\}_{i=1}^{p}$, the problem studied in this thesis is to construct $\widehat{w}:\R\times(0,\infty)\to(0,\infty)$ that (i) fits the quotes, (ii) satisfies the conditions of \cref{sec:noarb}, and (iii) generalises to quotes not used in the construction. Conditions (i)--(iii) conflict; quantifying the trade-off across method families is the object of the thesis.

\chapter{Related Work}\label{ch:literature}

\section{Parametric surfaces: SVI and SSVI}
The Stochastic Volatility Inspired (SVI) parameterisation \cite{gatheral2004parsimonious} models a total-variance slice with five parameters,
\begin{equation}\label{eq:sviraw}
  w(k) \;=\; a + b\left(\rho\,(k-m) + \sqrt{(k-m)^2+\sigma^2}\right),
\end{equation}
and owes its longevity to closed-form derivatives, trader-interpretable reparameterisations, and its exact large-maturity consistency with the Heston smile. Gatheral and Jacquier \cite{gatheral2014arbitrage} supply the modern no-arbitrage theory: explicit butterfly conditions via the Durrleman function, three equivalent parameterisations (raw, natural, Jump--Wings) with exact conversions, an algorithm that \emph{eliminates} butterfly arbitrage by adjusting the Jump--Wings call slope and minimum variance, and the surface-level SSVI family
\begin{equation}\label{eq:ssvi}
  w(k,\theta_\tau) \;=\; \frac{\theta_\tau}{2}\Big(1+\rho\varphi(\theta_\tau)k+\sqrt{(\varphi(\theta_\tau)k+\rho)^2+1-\rho^2}\Big),
\end{equation}
where $\theta_\tau$ is the at-the-money total variance and $\varphi$ a decreasing function, with sharp theorems characterising when \cref{eq:ssvi} is entirely free of static arbitrage. Calibration is delicate---the SVI objective has spurious local minima---and the quasi-explicit reformulation of \cite{zeliade2009quasi} restores convexity of the inner subproblem. Chapter~\ref{ch:parametric} replicates this toolkit in full.

\section{Per-surface neural smoothers}
Ackerer, Tagasovska and Vatter \cite{ackerer2020deep} introduced the architecture adopted in Chapter~\ref{ch:deepsmoother}: the implied-volatility (or total-variance) surface is modelled as the \emph{product of a prior and a neural corrector}, and the no-arbitrage conditions are imposed as \emph{soft constraints}---penalties on the network's derivatives evaluated on a grid wider than the quotes. Chataigner, Cr\'epey and Dixon \cite{chataigner2020deep} work in price space, penalising the Dupire-relevant derivatives $(\partial^2_{KK}P\ge0$, $\partial_T P\ge0)$ and extracting local volatility by automatic differentiation; importantly they document that \emph{hard}-constrained architectures (e.g.\ splines of non-negative weights) lose too much expressiveness, motivating the soft-constraint consensus. Hoshisashi, Phelan and Barucca \cite{hoshisashi2024noarbitrage} systematise the approach (DCNN): exact first- and second-order derivatives via automatic differentiation, hinge penalties for each no-arbitrage inequality evaluated on a collocation set distinct from the data, and twice-differentiable activations---ReLU being excluded since its second derivative vanishes almost everywhere.

\section{Neural operators}
Neural operators \cite{kovachki2023neural} learn maps between function spaces with the defining property of \emph{discretisation invariance}: the same trained model accepts inputs sampled on arbitrary grids. Operator deep smoothing \cite{gonon2024operator} applies a graph neural operator (GNO) to implied-volatility smoothing: each day's quotes are treated as a pointwise discretisation of a latent surface, and a single $\sim$100k-parameter model, trained on roughly a decade of 20-minute CBOE S\&P~500 snapshots ($\sim$60 million quotes), smooths any quote set in one forward pass. The architecture modifies the standard GNO by making the first layer purely non-local, which allows the hidden state to be \emph{created} at arbitrary output locations (interpolation); training uses a vega-weighted relative fitting loss and finite-difference no-arbitrage penalties. Chapter~\ref{ch:operator} adapts this approach to the daily-data regime available to academic research.

\section{Positioning}
The three families have never been compared on equal terms: \cite{gatheral2014arbitrage} predates the neural literature, \cite{ackerer2020deep} and \cite{hoshisashi2024noarbitrage} benchmark against their own SVI implementations on different datasets, and \cite{gonon2024operator} compares to Ackerer's \emph{reported} numbers rather than a shared sample. This thesis contributes the missing controlled comparison, together with an arbitrage audit of uniform stringency and the low-data transfer study of Chapter~\ref{ch:operator}.

\chapter{Data}\label{ch:data}

\section{Sources}
All market data are from OptionMetrics IvyDB US, accessed through WRDS. Four tables are used: daily option quotes on the S\&P~500 index (SPX), forward prices, the zero-coupon yield curve, and---as an industry reference only---OptionMetrics' own smoothed volatility surface. The sample covers 2018--2022: a calm regime (2019), the COVID-19 crash (March 2020), the recovery, and the 2022 sell-off, enabling the per-regime analysis of Chapter~\ref{ch:evaluation}.

\section{Cleaning pipeline}\label{sec:pipeline}
The raw exports are processed with a lazy, streaming Polars pipeline (Notebook~01) whose main steps are:
\begin{enumerate}
  \item \textbf{Typing and units.} OptionMetrics stores strikes in thousandths (a displayed strike of $4500$ is stored as $4500000$); dates are parsed defensively; quotes with $\mathrm{bid}\le 0$, $\mathrm{ask}<\mathrm{bid}$, or maturities outside $[7, 730]$ days are discarded.
  \item \textbf{Missing implied volatilities.} The fields \texttt{impl\_volatility}, \texttt{delta} and \texttt{vega} are missing on exactly the same rows ($\sim$5.7\%); $99.98\%$ of these are deep in-the-money options for which OptionMetrics does not attempt the inversion (the time value is buried in the spread). Since the surface is built from out-of-the-money quotes, these rows are dropped; the 118 OTM exceptions (out of 8.5 million rows) are immaterial.
  \item \textbf{Contract types.} Standard SPX (third-Friday, AM-settled) and SPXW weeklys (PM-settled) may share an expiration date with slightly different prices. The \texttt{root} field being unpopulated in the export, contracts are labelled by the settlement flag, and third-Friday collisions are resolved by preferring the standard contract; weeklys are retained on all other expirations, densifying the short end. One expiration therefore maps to exactly one contract.
  \item \textbf{Coordinates and discounting.} Forwards are taken from the quote file and completed by a join with the forward table; $k=\log(K/F)$; $\tau$ in ACT/365; the discount factor $D=e^{-r\tau}$ uses the zero curve at the nearest tenor.
  \item \textbf{Validation.} Implied volatilities are recomputed independently by a vectorised Black-inversion for every retained quote and compared to OptionMetrics' values: the median absolute gap is below $10^{-3}$ in volatility, validating forwards, discounting and coordinates simultaneously.
\end{enumerate}
The output is a typed Parquet dataset (one row per quote: date, expiry, $\tau$, $k$, bid/ask/mid, spread, volume, open interest, implied volatility, discount) on which all subsequent chapters operate; slices are always formed by exact expiration date, never by rounded maturity.

% FIGURE (export from NB01): quote density over time
\begin{figure}[t]
  \centering
  \maybegraphics[width=0.8\textwidth]{figures/nb01_quote_density.png}
  \caption[OTM quote density]{Number of out-of-the-money SPX quotes per trading day after cleaning. The weekly sawtooth reflects expiration cycles.}
  \label{fig:density}
\end{figure}

\section{Evaluation domains}\label{sec:domains}
Two evaluation domains are carried through the thesis. The \emph{full} domain retains every cleaned quote and is honest about real difficulty; the \emph{paper} domain restricts to $k\in[-0.6,0.4]$ and $\mathrm{DTE}\ge 10$, matching the effective domains of the reference papers and enabling literature-comparable numbers. Results are reported on both.

% =====================================================================
%  Chapter: The Parametric Benchmark — SVI and SSVI
% =====================================================================
\chapter{The Parametric Benchmark: SVI and SSVI}
\label{ch:parametric}

Before any neural method can claim an improvement, the parametric state of the
art must be implemented \emph{properly} and audited \emph{honestly}. This
chapter replicates the complete toolkit of \cite{gatheral2014arbitrage} and
turns it into a quantified benchmark on the real SPX dataset of
\cref{ch:data}. Every replicated component is validated against the numbers
printed in the original paper.

\section{Replication of Gatheral--Jacquier (2014)}
\label{sec:replication}

\subsection{Parameterizations and exact conversions}
We implement the three equivalent SVI parameterizations: \emph{raw}
(\cref{eq:sviraw}), \emph{natural}, and \emph{jump-wings} (JW), together with
the closed-form conversions of Lemmas~3.1--3.2 of the paper. The JW
parameters $(v_t,\psi_t,p_t,c_t,\tilde v_t)$ --- ATM variance, ATM skew, wing
slopes and minimum variance --- are the trader-facing coordinates and the
vehicle of the butterfly-elimination algorithm below.

\begin{remark}[A correction to the printed Lemma 3.2]
\label{rem:lemma32}
When inverting the JW map back to raw parameters we found that the printed
formula omits one term; using it verbatim fails to recover the raw parameters
of the paper's own example. \Cref{app:lemma32} derives the corrected inverse,
which handles the case $m<0$ through the identity
$\sqrt{m^2+\sigma^2}=m/\beta$. All conversions were validated by round-trips
over $2{,}000$ random parameter draws, with maximum error
$2.7\times10^{-10}$, and the JW values of the paper's Vogt example
$(0.01742625,\,-0.1752111,\,0.6997381,\,1.316798,\,0.0116249)$ are reproduced
to the printed digit.
\end{remark}

\subsection{Static-arbitrage diagnostics and the Vogt smile}
The butterfly criterion is the Durrleman condition $g(k)\ge0$ of
\cref{eq:durrleman} together with the wing condition $b(1+|\rho|)<2$
(Lee's moment bound \cite{lee2004moment}); calendar arbitrage between two
slices is measured by their \emph{crossedness}. Applied to the Vogt smile
(Example~3.1 of the paper), the diagnostics reproduce the textbook pathology:
a perfectly innocent-looking total-variance slice whose $g$ dips to
$\min_k g=-0.0329$, i.e.\ a locally \emph{negative} implied density
(\cref{fig:vogt}, already shown in \cref{ch:foundations}).

\subsection{Quasi-explicit calibration}
Our production calibrator follows the quasi-explicit method of
\cite{zeliade2009quasi}: freezing $(m,\sigma)$ and substituting
$y=(k-m)/\sigma$ makes raw SVI \emph{linear} in $(a,d,c)=(a,\rho b\sigma,
b\sigma)$, so the inner problem is a convex least squares under linear
no-arbitrage constraints (solved by SLSQP with an analytic gradient), and only
$(m,\sigma)$ are searched by Nelder--Mead with restarts. On a synthetic SSVI
slice observed with $1\%$ multiplicative noise the calibrator recovers the
slice with $\mathrm{RMSE}(w)=2.8\times10^{-4}$ and $\min_k g=0.44>0$.

\subsection{The paper's recipe, butterfly elimination, SSVI variants,
interpolation}
Four further components complete the replication.
\emph{(i)} The \S5.2 calibration recipe --- square-root-SVI initialization
followed by slice-by-slice refinement under a heavy crossedness penalty ---
removes an artificially injected calendar crossing (total crossedness
$0.113\to0.000$; \cref{fig:recipe}).
\emph{(ii)} The \S5.1 butterfly-elimination algorithm in JW coordinates,
$c_t'=p_t+2\psi_t$ and $\tilde v_t'=v_t\,4p_tc_t'/(p_t+c_t')^2$, reproduces
Example~5.1 exactly: $(c_t',\tilde v_t')=(0.3493158,\,0.0154818)$, taking the
Vogt smile from $\min g=-0.033$ to $+0.270$ (\cref{fig:elimination}).
\emph{(iii)} The three SSVI variants --- power-law, Heston-like, and the
modified power-law of Eq.~(4.5), fully arbitrage-free whenever
$\eta(1+|\rho|)\le2$ --- are verified against Theorems~4.1--4.2.
\emph{(iv)} The arbitrage-free interpolation of Lemma~5.1 (price-space
interpolation with $\sqrt\theta$ weights) and the extrapolation of
Theorem~4.3 are implemented and audited ($\min g = 0.52$ on the interpolated
slice).

\begin{figure}[t]
  \centering
  \maybegraphics[width=\textwidth]{figures/nb02_butterfly_elimination.png}
  \caption[Butterfly elimination via jump-wings (Example 5.1)]{Butterfly
  elimination in JW coordinates, reproducing Example~5.1 of
  \cite{gatheral2014arbitrage}. Left: the original (red) and repaired (green)
  smiles are nearly indistinguishable; middle: $g(k)$ is lifted above zero;
  right: the implied density becomes non-negative.}
  \label{fig:elimination}
\end{figure}

\begin{figure}[t]
  \centering
  \maybegraphics[width=\textwidth]{figures/nb02_recipe_crossedness.png}
  \caption[The \S5.2 recipe: crossedness penalty]{The paper's calibration
  recipe. Without the penalty (left) an inflated short slice crosses its
  neighbours (calendar arbitrage); the crossedness penalty (right) removes
  the crossing at negligible fit cost.}
  \label{fig:recipe}
\end{figure}

\section{Benchmark methodology on real data}
\label{sec:benchmethod}

The benchmark driver applies the replicated machinery to
\texttt{option\_prices\_clean.parquet}, one calibration problem per trading
day, under the following protocol. Slices are formed by \emph{exact
expiration date} (\texttt{exdate}), never by rounded maturity --- the SPX/SPXW
resolution of \cref{ch:data} guarantees one contract family per slice.
Quotes are weighted by inverse bid--ask spread, the liquidity weighting used
throughout the reference papers. Within each slice $20\%$ of quotes are
withheld as a \emph{hold-out} set, so that the parametric benchmark is scored
on the same generalization footing as the learning-based methods of
\cref{ch:deepsmoother,ch:operator}.

For the surface-level SSVI fit, the ATM total variance $\theta_i$ of each
slice is estimated \emph{from the quotes} (interpolation at $k=0$), not from
the fitted SVI's extrapolation, which proved erratic for long maturities with
few near-the-money quotes. Monotonicity of the term structure
$\theta(\tau)$ is diagnosed \emph{in maturity order} --- an earlier version of
the pipeline sorted the $\theta$'s before testing monotonicity, silently
masking genuine inversions --- raw inversions are counted and reported, and
repaired by isotonic regression (pool-adjacent-violators) before the global
$(\rho,\eta,\gamma)$ fit under the Theorem~4.2 butterfly constraints.
Finally, the audit covers both arbitrage directions: the Durrleman $g$
within each slice, and the crossedness between adjacent fitted slices.

All results are reported on two evaluation domains: the \textbf{full} quoted
domain, and the \textbf{paper} domain $k\in[-0.6,0.4]$, $\mathrm{DTE}\ge10$
used by the reference literature. The former is honest about the real
difficulty of the data; the latter makes our numbers directly comparable to
published work.

\section{Results}
\label{sec:benchresults}

\TODO{The numbers below are from the preliminary 8-day January-2018 sample
(LIMIT\_DATES=8); replace with the full 2018--2022 run before submission.}

\begin{table}[t]
  \centering
  \caption[Parametric benchmark: fit and arbitrage audit]{Parametric
  benchmark on SPX (median IV RMSE in volatility points; preliminary 8-day
  sample, 258 slices). Hold-out $\approx$ in-sample confirms the parametric
  fits do not overfit.}
  \label{tab:bench}
  \begin{tabular}{lS[table-format=1.3]S[table-format=1.3]S[table-format=1.3]S[table-format=1.3]}
    \toprule
    & \multicolumn{2}{c}{Full domain} & \multicolumn{2}{c}{Paper domain} \\
    \cmidrule(lr){2-3}\cmidrule(lr){4-5}
    Method & {In-sample} & {Hold-out} & {In-sample} & {Hold-out} \\
    \midrule
    SVI (per slice)   & 0.949 & 0.968 & 0.806 & 0.806 \\
    SSVI (per day)    & 2.789 & 2.790 & 2.551 & 2.670 \\
    \midrule
    Butterfly-violating slices (\%) & \multicolumn{2}{c}{6.20} & \multicolumn{2}{c}{2.44} \\
    Days with raw-$\theta$ inversions (\%) & \multicolumn{2}{c}{12} & \multicolumn{2}{c}{12} \\
    Median cross-slice crossedness & \multicolumn{2}{c}{$3.3\times10^{-2}$} & \multicolumn{2}{c}{$1.1\times10^{-2}$} \\
    \bottomrule
  \end{tabular}
\end{table}

\Cref{tab:bench} summarizes the benchmark. Three findings deserve emphasis.

\paragraph{The fit--arbitrage trade-off is now a number.} Per-slice SVI
achieves a median hold-out error of $0.97$ vol points on the full domain but
violates the butterfly condition on $6.2\%$ of slices and exhibits non-zero
crossedness between adjacent slices --- per-slice calibration produces
\emph{no} surface-level coherence. SSVI, arbitrage-free across maturities by
construction, pays for its guarantee with an error almost three times larger.
The window between these two numbers ($\approx0.97$ vs $\approx2.8$ vol
points) is exactly the target the learning-based methods must enter.

\paragraph{Difficulty is localized --- and not where expected.}
\Cref{tab:buckets} breaks the full-domain results by maturity bucket. Slices
beyond 180 days are essentially solved ($0.13$ vol pts, no violations); the
butterfly violations concentrate massively in the ultra-short bucket
(7--14 days, $55\%$ of slices). Strikingly, localizing the minimizer of $g$
shows that \emph{all} violations sit in the core moneyness region
$k\in[-0.6,0.4]$, not in the wings: at these maturities the SPX smile has a
pronounced kink at the money, total variance is small, and the $1/w$ term in
\cref{eq:durrleman} amplifies any excess convexity of the fit. Restricting
the moneyness domain therefore does \emph{not} remove the problem; the
$\mathrm{DTE}\ge10$ filter of the paper domain does most of the cleaning.

\begin{table}[t]
  \centering
  \caption[Benchmark by maturity bucket]{Full-domain results by maturity
  bucket. Violations concentrate in ultra-short maturities and, at those
  maturities, in the \emph{core} moneyness region (16 of 16 located at
  $k\in[-0.6,0.4]$), pointing to a structural limitation of the SVI
  parametric form at the short end rather than a calibration failure.}
  \label{tab:buckets}
  \begin{tabular}{lS[table-format=1.3]S[table-format=2.1]S[table-format=3.0]}
    \toprule
    Maturity bucket & {Median RMSE (vol pts)} & {Violations (\%)} & {Slices} \\
    \midrule
    7--14 d   & 0.995 & 55.2 & 29 \\
    15--60 d  & 1.003 & 0.0  & 109 \\
    61--180 d & 0.881 & 0.0  & 56 \\
    $>$180 d  & 0.135 & 0.0  & 64 \\
    \bottomrule
  \end{tabular}
\end{table}

\begin{figure}[t]
  \centering
  \maybegraphics[width=\textwidth]{figures/nb02_buckets.png}
  \caption[Difficulty by maturity bucket]{Fit error (left) and
  butterfly-violation rate (right) by maturity bucket on the full domain.
  The arbitrage risk of the parametric benchmark is an ultra-short-maturity
  phenomenon.}
  \label{fig:buckets}
\end{figure}

\paragraph{The benchmark matches the industrial reference.} Evaluating our
per-day SVI fits at the grid points of the proprietary OptionMetrics
volatility surface (forwards interpolated from the zero-curve and forward
tables of \cref{ch:data}) yields a median absolute gap of $0.93$ vol points
over $2{,}958$ grid points ($90$th percentile $3.13$). A fully transparent,
replicable pipeline thus reproduces the commercial smoother to within one
volatility point at the median.

\begin{figure}[t]
  \centering
  \maybegraphics[width=0.85\textwidth]{figures/nb02_domains.png}
  \caption[Benchmark under both evaluation domains]{Median IV RMSE under the
  full and paper evaluation domains, in-sample and hold-out. The thesis
  reports both throughout.}
  \label{fig:domains}
\end{figure}

\begin{figure}[t]
  \centering
  \maybegraphics[width=\textwidth]{figures/nb02_day_smiles.png}
  \caption[A full trading day under the parametric benchmark]{One trading day:
  market quotes (dots), per-slice SVI (solid) and surface SSVI (dashed), with
  the day's ATM total-variance term structure. The raw $\theta(\tau)$
  estimates exhibit the inversions quantified in \cref{tab:bench}.}
  \label{fig:daysmiles}
\end{figure}

\section{Known limitation and outlook}
The SSVI objective is currently fitted by unweighted least squares while the
SVI slices use inverse-spread weights; this asymmetry slightly disadvantages
SSVI in \cref{tab:bench} and will be removed for the final full-sample run.
More fundamentally, the chapter's conclusion stands regardless: a single
$(\rho,\eta,\gamma)$ triple across $\sim$30 expirations is structurally too
rigid for SPX, while per-slice SVI buys accuracy at the price of measurable
static arbitrage. \Cref{ch:deepsmoother} asks whether a neural smoother can
occupy the space between.

% =====================================================================
%  Chapter: Derivative-Constrained Deep Smoothing
% =====================================================================
\chapter{Derivative-Constrained Deep Smoothing}
\label{ch:deepsmoother}

The parametric benchmark of \cref{ch:parametric} quantified a trade-off:
per-slice SVI fits well but leaks static arbitrage; SSVI is arbitrage-free by
construction but three times less accurate. This chapter implements a neural
smoother designed to occupy the space between, combining the
\emph{architecture} of \cite{ackerer2020deep} with the \emph{training
mechanics} of \cite{hoshisashi2024noarbitrage}.

\section{Method}
\label{sec:dsmethod}

\subsection{Architecture: an SSVI prior times a neural corrector}
Following \cite{ackerer2020deep}, the total-variance surface is modelled
multiplicatively,
\begin{equation}
  w_\theta(k,\tau) \;=\; w_{\mathrm{SSVI}}(k,\tau)\,\times\,
  \mathcal C_\theta(k,\tau),
  \qquad
  \mathcal C_\theta = \exp\!\big(\mathrm{MLP}_\theta(k,\tau)\big) > 0,
  \label{eq:priorcorrector}
\end{equation}
where the prior is an SSVI surface with a smooth power-law term structure
$\theta(\tau)=\alpha\tau^{\beta}$ fitted to the day's quotes, and the
corrector is a small multilayer perceptron ($2\times64$ hidden units,
$\tanh$ activations) whose inputs $(k,\tau)$ are rescaled to $[-1,1]^2$.
With standard small-weight initialization $\mathcal C_\theta\approx1$, so
training \emph{starts at the prior} and departs from it only where the data
demand. The choice of $\tanh$ is not cosmetic: the loss below involves
$\partial^2_k w_\theta$, and piecewise-linear activations such as ReLU have
vanishing second derivatives almost everywhere, which would silence the
butterfly penalty \cite{hoshisashi2024noarbitrage}.

\subsection{Loss: soft no-arbitrage penalties on a collocation grid}
Training minimizes
\begin{equation}
  \mathcal L(\theta)
  = \underbrace{\frac1N\sum_{i=1}^{N}\big(w_\theta(k_i,\tau_i)-w_i\big)^2}_{\text{fit on sparse quotes}}
  + \lambda_{\mathrm{bfly}}\,\overline{\relu\!\big(-g_\theta\big)^2}\Big|_{\hat X}
  + \lambda_{\mathrm{cal}}\,\overline{\relu\!\big(-\partial_\tau w_\theta\big)^2}\Big|_{\hat X},
  \label{eq:dsloss}
\end{equation}
where $g_\theta$ is the Durrleman function of \cref{eq:durrleman} evaluated
on the model surface, and --- the decisive detail imported from
\cite{hoshisashi2024noarbitrage} --- the penalties are averaged over a
\emph{dense collocation grid} $\hat X$ ($24\times10$ points spanning the
day's quoted rectangle) that is disjoint from the quotes. No-arbitrage is
thus enforced everywhere, including regions with no data, which is what makes
extrapolation safe. The derivatives $w_k,w_{kk},w_\tau$ entering
\cref{eq:dsloss} are computed by \emph{exact} automatic differentiation of
the network (nested \texttt{elementwise\_grad}); before any training we
verified the machinery against the closed-form SVI derivatives of
\cref{ch:foundations}, with agreement at the $10^{-16}$ level. Soft
penalties, rather than hard architectural constraints, follow the evidence of
\cite{chataigner2020deep} that hard constraints sacrifice too much
representational power; we set $\lambda_{\mathrm{bfly}}=\lambda_{\mathrm{cal}}=10$.

\section{Controlled experiments on synthetic surfaces}
\label{sec:dsablation}

The method is first validated where ground truth is known: a target SSVI
surface sampled on sparse irregular grids ($6$--$14$ quotes per slice, $1.5\%$
noise), with a deliberately misspecified prior fitted from the same sparse
quotes. Three ablations isolate the two design choices.

\begin{table}[t]
  \centering
  \caption[Deep smoother ablations]{Ablation study on a synthetic surface
  (IV RMSE against the dense ground truth, in vol points; $\min g$ and
  violation share on the collocation grid).}
  \label{tab:ablation}
  \begin{tabular}{lS[table-format=1.3]S[table-format=+1.3]S[table-format=2.1]}
    \toprule
    Variant & {RMSE (vol pts)} & {$\min g$} & {Violations (\%)} \\
    \midrule
    Full model (prior + constraints) & 0.187 & +0.33 & 0.0 \\
    No constraints ($\lambda=0$)     & 0.247 & {---} & {---} \\
    No prior (flat BS prior)         & 1.418 & {---} & {---} \\
    \bottomrule
  \end{tabular}
  \\[2pt]\footnotesize \TODO{fill the $\min g$/violation columns of the
  ablated variants from the notebook run.}
\end{table}

Two lessons emerge from \cref{tab:ablation}. First, the \emph{prior is the
dominant ingredient} in the sparse-data regime: without it the corrector must
invent the entire surface shape and the error grows sevenfold. Second, the
constraints act as a \emph{regularizer} as well as a safety device --- the
fully constrained model fits \emph{better} than the unconstrained one, in
line with the observations of \cite{hoshisashi2024noarbitrage}, while
guaranteeing a clean $g$ on the collocation grid
(\cref{fig:gheatmap}).

\begin{figure}[t]
  \centering
  \maybegraphics[width=\textwidth]{figures/nb03_gheatmap.png}
  \caption[Arbitrage heat-map of the deep smoother]{Durrleman $g(k,\tau)$
  heat-map. Left: the constrained model is non-negative everywhere. Right:
  removing the penalties lets butterfly arbitrage (red) leak into regions
  unsupported by data.}
  \label{fig:gheatmap}
\end{figure}

\begin{figure}[t]
  \centering
  \maybegraphics[width=0.8\textwidth]{figures/nb03_tradeoff.png}
  \caption[Training dynamics: fit vs penalty]{Training dynamics of the fit
  term and butterfly penalty (log scale), in the style of Fig.~7 of
  \cite{hoshisashi2024noarbitrage}.}
  \label{fig:tradeoff}
\end{figure}

\paragraph{Robustness to sparsity.} Subsampling the synthetic quotes exposes
the structural advantage of a \emph{surface-level} smoother over per-slice
calibration: at $25\%$ of the quotes, no slice retains the six points needed
to calibrate a five-parameter SVI --- per-slice calibration simply ceases to
exist --- while the deep smoother, which pools information across maturities,
still reconstructs the surface at $0.315$ vol points
(\cref{fig:sparsity}).

\begin{figure}[t]
  \centering
  \maybegraphics[width=0.7\textwidth]{figures/nb03_sparsity.png}
  \caption[Robustness to quote sparsity]{Reconstruction error as quotes are
  subsampled. Per-slice SVI loses calibrable slices and eventually cannot be
  evaluated; the surface-level smoother degrades gracefully.}
  \label{fig:sparsity}
\end{figure}

\section{Real-data protocol and preliminary results}
\label{sec:dsreal}

On real SPX days the protocol mirrors the benchmark exactly: OTM quotes
grouped by expiration, an SSVI prior fitted per day, $20\%$ of quotes
withheld, penalties on the day's collocation grid, and the same
arbitrage audit. On the preliminary sample the smoother attains a hold-out
error of $0.158$ vol points against $0.153$ in-sample --- hold-out
$\approx$ in-sample, as for the parametric fits, with $\min g>0$ on the
audit grid. \TODO{replace with full-sample per-day results joined against
\texttt{benchmark\_svi\_slices\_\{full,paper\}.parquet}, and add the
comparison figure deep-vs-SVI-vs-SSVI per day.}

\begin{figure}[t]
  \centering
  \maybegraphics[width=\textwidth]{figures/nb03_smiles.png}
  \caption[Deep smoother on real SPX quotes]{The deep smoother on a real
  trading day: quotes (dots) and fitted smiles per maturity.}
  \label{fig:dssmiles}
\end{figure}

\section{Discussion}
The prior-times-corrector design earns both of its components: the prior
carries the sparse-data regime, the derivative penalties buy near-exact
no-arbitrage at zero cost in accuracy --- on synthetic data they even improve
it. The remaining honest caveat, inherited from the soft-constraint
paradigm and documented by \cite{chataigner2020deep}, is that penalties
\emph{reduce} violations without \emph{abolishing} them on unseen regions;
we therefore always report the residual violation rate alongside the error.
The chapter's limitation is economic rather than statistical: one network
must be trained \emph{per day}. Removing that cost is the subject of
\cref{ch:operator}.

% =====================================================================
%  Chapter: Operator Deep Smoothing
% =====================================================================
\chapter{Operator Deep Smoothing: One Model for Every Surface}
\label{ch:operator}

Every method so far --- SVI, SSVI, and the deep smoother of
\cref{ch:deepsmoother} --- learns \emph{one surface at a time}: a new day
requires a new calibration or a new training run. Following the operator
deep smoothing approach of \cite{gonon2024operator}, this chapter learns the
\emph{smoothing map itself},
\begin{equation}
  F_\theta:\ \{(k_i,\tau_i,\sigma_i)\}_{i=1}^{p}\ \longmapsto\
  \hat v(\cdot,\cdot),
  \label{eq:operator}
\end{equation}
trained once across many days and applied to unseen days \emph{without
retraining}. The enabling property is \emph{discretization invariance}: the
operator accepts inputs of arbitrary size $p$ and arbitrary spatial layout
--- exactly the nature of daily option quotes.

\section{Two architectures}
\label{sec:oparch}

\subsection{An executable DeepONet-class operator}
Our in-notebook operator is of the DeepONet/DeepSets family
\cite{lu2021deeponet}: a permutation-invariant set encoder embeds each quote
$(k_i,\tau_i,\sigma_i)$ and \emph{mean-pools} the embeddings --- the same
averaging philosophy that underlies the Nystr\"om kernel sums of
\cite{gonon2024operator}, and the mechanism that renders the representation
invariant to input size and ordering --- while a trunk network embeds the
query coordinate $(k,\tau)$, so the surface can be evaluated anywhere. A
Softplus output enforces positivity. This model is deliberately lighter than
the graph operator below: it isolates the \emph{operator questions}
(no retraining, subsampling robustness, transfer) from the
\emph{architecture question} (graph kernels).

\subsection{The modified graph neural operator of the paper}
We additionally provide a faithful PyTorch implementation of the modified
GNO of \cite{gonon2024operator}: a purely \emph{non-local first layer}
(dropping the pointwise linear map, so hidden state can be \emph{created} at
arbitrary output locations --- the interpolation property), in-neighborhoods
drawn from input locations with truncation along the maturity axis and a
low-discrepancy subsampling cap of $K=50$ nodes, per-layer lifting and
projection networks, kernel networks with input skip connections, GELU
activations and Softplus output; at $J=4$ layers and $16$ channels the model
has $\approx10^5$ parameters, the paper's configuration. Training this
architecture at scale requires GPU time and is left for the final
full-sample experiments. \TODO{if GPU runs are completed before submission,
report GNO results alongside the DeepONet results below.}

\section{The paper's loss, resolved by finite differences}
Training minimizes the five-term loss of \cite{gonon2024operator}: a
\emph{vega-weighted relative} fitting term (deep-OTM quotes with negligible
vega must not dominate), the butterfly hinge
$\overline{(g-\varepsilon)^-}$ at $\varepsilon=10^{-3}$ with $g$ resolved by
central finite differences on a rectilinear grid, a calendar penalty
(monotone total variance across maturities --- the scale-free ratio form of
the paper is its equivalent in transformed coordinates), and two curvature
regularizers $\|\partial^2_\tau\hat v\|_2,\|\partial^2_k\hat v\|_2$, with the
paper's weights $(1,10,10,0.01,0.01)$. The finite-difference Durrleman
machinery is validated against the closed forms of \cref{ch:foundations}
before use. During training every surface is randomly subsampled to
$60$--$100\%$ of its quotes --- the augmentation that forges robustness to
sparse inputs.

\section{The low-data transfer experiment}
\label{sec:transfer}

\cite{gonon2024operator} train on roughly $6\times10^{7}$ intraday CBOE
points spanning a decade. Academic access (OptionMetrics) is \emph{daily}:
about $1{,}250$ surfaces for 2018--2022 --- two orders of magnitude less.
This gap is the research question of the chapter:

\begin{quote}\itshape
Does operator deep smoothing work in the low-data daily regime, and does
pre-training on synthetic arbitrage-free surfaces close the gap?
\end{quote}

The experiment crosses three data regimes under identical architecture,
loss and budget (\cref{tab:regimes}), with a \emph{strict chronological
split} --- training on the earliest dates, testing on the latest --- so that
the test period contains a genuine volatility-regime shift, as it does in
reality.

\begin{table}[t]
  \centering
  \caption[Transfer-experiment regimes]{The three training regimes of the
  transfer experiment. The synthetic bank consists of SSVI surfaces with
  parameters drawn from realistic ranges, sampled on sparse irregular grids
  with multiplicative noise --- arbitrage-free teachers in unlimited supply.}
  \label{tab:regimes}
  \begin{tabular}{llll}
    \toprule
    Regime & Pre-training & Fine-tuning & Question \\
    \midrule
    R1 & --- & real train days & operator at academic data scale \\
    R2 & synthetic bank & --- & zero-shot synthetic$\to$real transfer \\
    R3 & synthetic bank & real train days & does pre-training close the gap? \\
    \bottomrule
  \end{tabular}
\end{table}

\begin{table}[t]
  \centering
  \caption[Transfer-experiment results]{Transfer results on unseen
  chronological test days (IV RMSE in vol points; arbitrage audit on the
  dense grid). \TODO{fill from the full-sample notebook run.}}
  \label{tab:transfer}
  \begin{tabular}{lS[table-format=1.3]S[table-format=+1.3]S[table-format=2.1]}
    \toprule
    Regime & {IV RMSE} & {$\min g$} & {Violations (\%)} \\
    \midrule
    R1 real-only            & {\TODO{}} & {\TODO{}} & {\TODO{}} \\
    R2 synth-only (0-shot)  & {\TODO{}} & {\TODO{}} & {\TODO{}} \\
    R3 pretrain + finetune  & {\TODO{}} & {\TODO{}} & {\TODO{}} \\
    \bottomrule
  \end{tabular}
\end{table}

\begin{figure}[t]
  \centering
  \maybegraphics[width=0.75\textwidth]{figures/nb04_transfer.png}
  \caption[Transfer experiment R1/R2/R3]{IV RMSE on unseen test days under
  the three training regimes. Whatever the ordering, the result answers a
  question the original paper does not address; the expected pattern under a
  regime shift is $\mathrm{R3}\le\mathrm{R1}<\mathrm{R2}$.}
  \label{fig:transfer}
\end{figure}

\section{Discretization invariance}
\label{sec:invariance}

The operator's signature property is tested directly: one unseen test day is
fed $100\%$, $50\%$ and $25\%$ of its quotes, and the movement of the
\emph{predicted dense surface} is measured. A per-day method cannot even
enter this comparison --- each subsample would require a fresh training run.
\Cref{fig:invariance} reports the mean surface shift; the operator's output
remains stable as the input thins, the empirical counterpart of the
discretization-invariance property of \cite{kovachki2023neural}.

\begin{figure}[t]
  \centering
  \maybegraphics[width=\textwidth]{figures/nb04_invariance.png}
  \caption[Discretization invariance]{Left: mean shift of the predicted
  surface as the input is subsampled. Right: a smoothed slice with full
  versus quarter input.}
  \label{fig:invariance}
\end{figure}

\begin{figure}[t]
  \centering
  \maybegraphics[width=0.8\textwidth]{figures/nb04_surface.png}
  \caption[Operator-smoothed surface on an unseen day]{Operator-smoothed
  implied-volatility surface on an unseen test day --- produced by a single
  forward pass, with no day-specific training.}
  \label{fig:opsurface}
\end{figure}

\section{Discussion}
The operator changes the economics of smoothing: the cost of a new surface
drops from a calibration (\cref{ch:parametric}) or a training run
(\cref{ch:deepsmoother}) to a single forward pass, and the same weights serve
every day. The price is a heavier training pipeline, the loss of per-day
interpretability, and --- in the academic data regime --- the exposure to
overfitting that the transfer experiment is designed to measure. The
unified comparison of all methods, including the proprietary OptionMetrics
surface, is the subject of \cref{ch:evaluation}.

% =====================================================================
%  Chapter: Unified Evaluation
% =====================================================================
\chapter{Unified Evaluation}
\label{ch:evaluation}

The value of this thesis lies less in any single method than in their
comparison \emph{on equal terms}: the same SPX days, the same hold-out
protocol, the same arbitrage audit. This chapter assembles the per-day
outputs of the preceding chapters
(\texttt{benchmark\_svi\_slices\_\{full,paper\}.parquet},
\texttt{benchmark\_ssvi\_days\_\{full,paper\}.parquet},
\texttt{deep\_smoother\_days.parquet}, \texttt{operator\_days.parquet}) into
a single evaluation.

\section{Protocol}
Five surfaces are compared on every evaluation day: per-slice \textbf{SVI},
surface \textbf{SSVI}, the proprietary \textbf{OptionMetrics} smoothed
surface, the per-day \textbf{deep smoother} of \cref{ch:deepsmoother}, and
the \textbf{neural operator} of \cref{ch:operator} (evaluated only on its
chronological test window, where it has never seen the day). Metrics:
median IV RMSE on held-out quotes (vol points); butterfly-violation rate and
minimum Durrleman $g$ on a dense audit grid; cross-maturity crossedness;
and the computational cost of producing a new surface.

\section{Headline comparison}

\begin{table}[t]
  \centering
  \caption[Unified five-method comparison]{Unified comparison (median over
  evaluation days; hold-out IV RMSE in vol points). Preliminary values from
  the 8-day sample where available; \TODO{fill remaining cells from the
  full-sample runs of NB02--NB05 before submission.}}
  \label{tab:unified}
  \begin{tabular}{lS[table-format=1.3]S[table-format=2.1]l}
    \toprule
    Method & {Hold-out RMSE} & {Bfly viol.\ (\%)} & Cost of a new surface \\
    \midrule
    SVI (per slice)        & 0.968 & 6.2 & one calibration per slice \\
    SSVI (per day)         & 2.790 & 0.0 & one 3-parameter fit per day \\
    OptionMetrics surface  & {---} & {\TODO{}} & proprietary \\
    Deep smoother (per day)& {\TODO{}} & {\TODO{}} & one training run per day \\
    Neural operator        & {\TODO{}} & {\TODO{}} & one forward pass \\
    \bottomrule
  \end{tabular}
\end{table}

Interpretation of the headline table follows the arc of the thesis. The gap
between SVI and SSVI ($\approx0.97$ vs $\approx2.8$ vol points hold-out) is
the price of built-in no-arbitrage in the parametric world. The deep
smoother is designed to close most of that gap while keeping violations near
zero; the operator is designed to do so \emph{at negligible marginal cost per
day}. Against the industrial reference, our transparent SVI pipeline already
sits within $0.93$ vol points (median) of the OptionMetrics surface over
$2{,}958$ grid points.

\section{Analysis by market regime}
The 2018--2022 sample spans three qualitatively different regimes: the calm
of 2019, the COVID-19 crash of March 2020, and the 2022 rate-hike sell-off.
All per-day metrics are re-aggregated within each regime.
\TODO{insert regime table and figure from NB05 once the full-sample runs
complete; the working hypothesis, supported by the bucket analysis of
\cref{tab:buckets}, is that method rankings are stable but absolute errors
and violation rates degrade sharply in March 2020, and that the operator's
test window (2022) constitutes a genuine out-of-regime test.}

\section{Robustness to sparsity}
The synthetic sparsity study of \cref{sec:dsablation} is repeated on real
days by subsampling quotes before evaluation. Per-slice SVI degrades and
eventually loses calibrable slices; the surface-level methods degrade
gracefully; the operator additionally demonstrates stability of its output
under input subsampling (\cref{sec:invariance}) --- a property no per-surface
method possesses.


% =====================================================================
%  Chapter: Conclusion
% =====================================================================
\chapter{Conclusion}
\label{ch:conclusion}

\section{Summary of findings}
This thesis built a single, reproducible SPX pipeline --- from raw
OptionMetrics exports to a unified five-method evaluation --- and used it to
traverse the methodological spectrum of implied-volatility surface
construction.

On the parametric side, we replicated the complete toolkit of
\cite{gatheral2014arbitrage} to the digit, correcting in passing a term in
the printed inverse of Lemma~3.2, and turned it into an audited benchmark.
The benchmark quantifies the fit--arbitrage trade-off (SVI $\approx0.97$ vol
points hold-out with $6.2\%$ butterfly-violating slices; SSVI arbitrage-free
at $\approx2.8$ vol points) and localizes the difficulty: violations
concentrate in ultra-short maturities and, contrary to intuition, in the
\emph{core} moneyness region rather than the wings --- a structural
limitation of the SVI form where the ATM kink meets small total variance.
The transparent pipeline matches the proprietary OptionMetrics surface to
$0.93$ vol points at the median.

On the learning side, the derivative-constrained deep smoother --- an SSVI
prior multiplied by a positive neural corrector, trained with
exact-autodiff Durrleman and calendar penalties on a collocation grid ---
showed in controlled ablations that the prior dominates in the sparse-data
regime (RMSE $0.187$ vs $1.418$ without it) and that the constraints act as
a regularizer (beating the unconstrained variant, $0.187$ vs $0.247$) while
keeping $g\ge0$. Where per-slice SVI ceases to exist at $25\%$ quote
coverage, the surface-level smoother still reconstructs at $0.315$ vol
points. Finally, the neural operator lifts the remaining economic
constraint --- one training per day --- and its low-data transfer experiment
(real-only, synthetic-zero-shot, pretrain-plus-finetune) addresses a regime
the original paper does not: smoothing with two orders of magnitude less
data than the industrial setting.

\section{Limitations}
Three limitations bound the claims. First, several headline numbers in this
draft come from a preliminary 8-day January-2018 sample; the full 2018--2022
runs replace them before submission. Second, soft no-arbitrage constraints
\emph{reduce} violations but do not \emph{abolish} them outside the training
support \cite{chataigner2020deep}; we report residual rates rather than
claiming guarantees. Third, a small methodological asymmetry (unweighted
SSVI objective versus spread-weighted SVI) slightly disadvantages SSVI in
the current tables and is scheduled for removal.

\section{Future work}
Natural continuations include: training the full graph neural operator of
\cite{gonon2024operator} at GPU scale on the daily sample; hybrid hard--soft
constraint schemes; extending the smoother to price space with a Dupire
local-volatility extraction and half-variance regularization as in
\cite{chataigner2020deep}; uncertainty quantification via operator
ensembles; and intraday data, which would test whether the low-data
conclusions of the transfer experiment persist as the data regime approaches
the industrial one.



% =====================  APPENDIX  =====================
\appendix
% =====================================================================
%  Appendix
% =====================================================================
\chapter{Technical Appendix}
\label{app:technical}

\section{The corrected inverse of Lemma 3.2}
\label{app:lemma32}
Given JW parameters $(v_t,\psi_t,p_t,c_t,\tilde v_t)$ at maturity $t$, set
$w_t=v_tt$, $b=\tfrac{\sqrt{w_t}}{2}(c_t+p_t)$ and
$\rho=(c_t-p_t)/(c_t+p_t)$. Define
$\beta=\rho-2\psi_t\sqrt{w_t}/b\in(-1,1)$ and
$\alpha=\operatorname{sign}(\beta)\sqrt{1/\beta^2-1}$, so that
$\sigma=\alpha m$ and $\sqrt{m^2+\sigma^2}=m/\beta$ for either sign of $m$.
Substituting into the raw-SVI expressions for $v_t$ and $\tilde v_t$ yields
\begin{equation}
  (v_t-\tilde v_t)\,t
  \;=\; b\,m\left(\frac{1}{\beta}-\rho-\alpha\sqrt{1-\rho^2}\right),
  \label{eq:lemma32}
\end{equation}
which is solved for $m$ (the printed statement omits the
$-\alpha\sqrt{1-\rho^2}$ term inside the bracket); finally
$a=\tilde v_t t-b\sigma\sqrt{1-\rho^2}$. The correction was validated by
round-trips over $2{,}000$ random draws (maximum error $2.7\times10^{-10}$)
and by exact recovery of the raw Vogt parameters from their JW image.

\section{Quasi-explicit calibration: inner problem constraints}
\label{app:zeliade}
With $y=(k-m)/\sigma$ and $(a,d,c)=(a,\rho b\sigma,b\sigma)$ the inner
problem is
$\min_{(a,d,c)}\sum_i \omega_i\,\big(a+d\,y_i+c\sqrt{y_i^2+1}-w_i\big)^2$
subject to $0\le c\le4\sigma$, $|d|\le c$, $|d|\le4\sigma-c$ and
$0\le a\le\max_i w_i$ --- a convex quadratic program solved by SLSQP with
analytic gradient; the outer $(m,\log\sigma)$ search uses Nelder--Mead with
three restarts.

\section{Isotonic repair of the ATM term structure}
\label{app:pava}
Raw ATM total variances $\hat\theta_i$ estimated from quotes occasionally
violate monotonicity in maturity (12\% of sample days). Before the SSVI fit
they are projected onto the monotone cone by pool-adjacent-violators:
repeatedly merge adjacent blocks whose means are decreasing, replacing them
by their weighted mean, until the sequence is non-decreasing. The raw
inversion counts are reported as a diagnostic in \cref{tab:bench}.

\section{Reproducibility}
\label{app:repro}
The pipeline is a sequence of five notebooks --- NB01 data cleaning
(Polars, lazy streaming), NB02 parametric benchmark, NB03 deep smoother,
NB04 neural operator, NB05 unified evaluation --- exchanging Parquet files
with fixed schemas. All stochastic steps are seeded; hold-out splits are
derived from per-date seeds; chronological splits are used for all
learning-based methods. \TODO{add repository URL / archive reference and
package versions.}


% =====================  BIBLIOGRAPHY  =====================
\bibliographystyle{abbrv}
\bibliography{bibliography}

\end{document}
